\pdfvariable trailerid {[<C3542026090300000000000000000000><C3542026090300000000000000000000>]}
\documentclass[10pt]{article}
\usepackage[a4paper,margin=18mm]{geometry}
\usepackage{amsmath,amsthm,unicode-math,microtype,xurl}
\usepackage[hidelinks]{hyperref}
\providecommand{\CRevisionRound}{2}
\newtheorem{theorem}{Theorem}
\newcommand{\R}{\mathbb R}
\newcommand{\Q}{\mathbb Q}
\newcommand{\sn}{\operatorname{sn}}
\newcommand{\SO}{\operatorname{SO}}
\title{\vspace{-1.2em}The Lagrange Top: Cubic Nutation, Two-Phase Reconstruction, and Exact Closure}
\author{Source-local theorem reconstruction}
\date{2026-09-03}
\begin{document}
\maketitle
\begin{abstract}
For the positive-inertia Lagrange symmetric heavy top we close the regular
reduction--reconstruction problem.  The inclination is inverted by one Jacobi
function, both cyclic phases are complete third-kind elliptic integrals, and
the full configuration closes exactly when two normalized phase increments
are rational.  Pole and repeated-root degenerations are kept outside singular
Euler formulas.  Exact computation certifies conventions, not the continuum
proof.
\end{abstract}

\section{Frozen convention and reduced cubic}
Let \(A=I_1=I_2>0\), \(C=I_3>0\), and \(\gamma>0\).  In regular
\(z\)-\(y\)-\(z\) Euler angles put \(u=\cos\theta\), \(L=p_\phi\), and
\(G=p_\psi\).  Our convention is
\begin{equation}
 H=\frac{p_\theta^2}{2A}+\frac{(L-Gu)^2}{2A(1-u^2)}
   +\frac{G^2}{2C}+\gamma u.                                    \label{eq:H}
\end{equation}
The two momenta and \(E=H\) are conserved.  Since
\(p_\theta=A\dot\theta\), energy gives the reduced cubic owner
\begin{equation}
 A^2\dot u^2=P(u):=
 2A\left(E-\frac{G^2}{2C}-\gamma u\right)(1-u^2)-(L-Gu)^2.       \label{eq:P}
\end{equation}
In particular
\begin{equation}
P(1)=-(L-G)^2,\qquad P(-1)=-(L+G)^2.                            \label{eq:poles}
\end{equation}
The smooth Hamiltonian flow is complete: its potential is bounded on compact
\(\SO(3)\), and positive kinetic energy makes each energy sublevel compact.

\begin{theorem}[regular reduction]\label{thm:main}
A nonsteady regular inclination has finite period if and only if its allowed
component of \(P\geq0\) is bounded by two simple roots in \((-1,1)\).
In the pole-incompatible chamber \(L\ne\pm G\), label the roots
\[
 -1<r_1<r_2<1<r_3,\qquad
 P=2A\gamma(u-r_1)(u-r_2)(u-r_3).
\]
This is the complete root classification for the regular reduced motion.
\ifnum\CRevisionRound>0
The formulas below give its least reduced period, reconstruction, and full
closure test.
\fi
\end{theorem}

\paragraph{Proof of the reduction.}
Allowed positions are exactly \(P\geq0\).  At a simple turning root the time
integral is locally \(\int s^{-1/2}ds\), while a double root gives
\(\int s^{-1}ds\).  Equation~\eqref{eq:poles}, the positive leading
coefficient, and the cubic sign pattern give the stated ordering.  This also
separates finite nutation from a separatrix.

\ifnum\CRevisionRound>0
\section{Elliptic inversion and two phase increments}
Set
\[
 d=r_2-r_1,\quad R=r_3-r_1,\quad
 k^2=\frac{d}{R},\quad \nu^2=\frac{\gamma R}{2A},\quad
 Q_0=\sqrt{2A\gamma R}.
\]
Direct substitution proves
\begin{equation}
 u(t)=r_1+d\,\sn^2(\nu(t-t_0),k),\qquad
 T=\frac{4A}{Q_0}K(k).                                          \label{eq:sn}
\end{equation}
Indeed \(z=(u-r_1)/d\) obeys
\(A^2d^2\dot z^2=2A\gamma d^2Rz(1-z)(1-k^2z)\).

Hamilton's equations are
\[
\dot\phi=\frac{L-Gu}{A(1-u^2)},\qquad
\dot\psi=\frac{G}{C}-u\dot\phi .
\]
Define the finite complete integrals
\[
 I_N=\frac{2}{Q_0(1-r_1)}
 \Pi\left(\frac{d}{1-r_1},k\right),\qquad
 I_S=\frac{2}{Q_0(1+r_1)}
 \Pi\left(-\frac{d}{1+r_1},k\right).
\]
Partial fractions on the outbound and return legs give the two phase
reconstruction
\begin{align}
 \Delta\phi&=(L-G)I_N+(L+G)I_S,                                 \label{eq:dphi}\\
 \Delta\psi&=G(1/C-1/A)T+(G-L)I_N+(G+L)I_S.                    \label{eq:dpsi}
\end{align}

\begin{theorem}[full regular closure]\label{thm:closure}
The full phase state in \(T^*\SO(3)\) closes if and only if
\[
 \frac{\Delta\phi}{2\pi}\in\Q,\qquad
 \frac{\Delta\psi}{2\pi}\in\Q.
\]
If their reduced denominators are \(q_\phi,q_\psi\), the least number of
nutations in a full return is \(\operatorname{lcm}(q_\phi,q_\psi)\).
\end{theorem}
\begin{proof}
After one least return of \((u,\dot u)\), autonomy makes the two angle
increments constant.  Rational increments therefore suffice.  Conversely, a
full return must occur after an integer number of reduced returns.  For
\(0<\theta<\pi\), the \(z\)-\(y\)-\(z\) decomposition is unique modulo
independent \(2\pi\) shifts of \(\phi,\psi\), so both increments must be
rational multiples of \(2\pi\).
\end{proof}
\fi

\ifnum\CRevisionRound>1
\section{Degenerate faces and quantum boundary}
If \(P(u_0)=P'(u_0)=0\), initial data at \(u_0\) give steady precession with two
constant angular rates; its regular closure is their rational-ratio test.
A nonconstant branch approaching the same double root takes infinite time.
A triple root is another critical face and is not inserted into
\eqref{eq:sn}.

Equation~\eqref{eq:poles} makes \(L=G\) necessary at the north pole and
\(L=-G\) necessary at the south pole.  Euler angles are singular there, so
existence and uniqueness come from the smooth group Hamiltonian.  If \(L=G\)
and two turning roots remain interior, formulas
\eqref{eq:sn}--\eqref{eq:dpsi} have the finite limiting root \(r_3=1\);
an orbit that reaches the pole is reconstructed in a group chart, not by
dividing by \(1-u^2\).  Sleeping tops, \(G=0\), \(\gamma=0\) (the quadratic
free-top boundary), and \(A=C\) are separate faces.

The positive rigid-body metric yields a symmetric elliptic kinetic operator
on \(C^\infty(\SO(3))\).  Elliptic completeness on the closed manifold makes
it essentially self-adjoint; bounded real multiplication by \(\gamma u\)
preserves that property.  Its unique closure equals the Friedrichs operator
and has compact resolvent.  We claim no closed spectrum.

\section{Evidence, collision boundary, and scope}
The executable artifact checks twelve rational parameter rows, sixty
reconstruction probes, exact Sturm isolators, four Jacobi substitutions, and
the steady/pole faces.  A producer-independent checker and a symbolic lane
own every stored field; repaired-hash mutations test the schema.  These are
finite convention receipts, not an enumeration proof.

C186 is the gravity-free Euler top; C244 has no body-axis spin momentum; C344
has a different complex-amplitude Poisson owner; C349 is a holonomic sphere
oscillator.  Thus this theorem owns the rigid-body two-phase reconstruction.
The Route-A tuple is
\[
(\mathrm{A0\_FAIL},\mathrm{A1\_WEAK},\mathrm{A2\_FAIL},
 \mathrm{A3\_FAIL},\mathrm{A4\_NATURAL\_QUANTIZATION}),
\]
and the overall verdict is rejection.  Scope is
\texttt{NO\_BAD\_EULER\_OR\_ROOT\_NUMBER}; Route B is false.  No target local
data, Euler factor, root number, automorphy, target divisor or functional
equation, target-zero match, or Hilbert--P\'olya operator is asserted.

\section{Sources}
Cushman and Bates give an authoritative global Lagrange-top and reduction
lineage~\cite{CB}; Audin supplies an authoritative spinning-top integrability
source~\cite{Audin}.  The formulas and software here are independently
reconstructed, without a priority claim.
\begin{thebibliography}{9}
\small
\bibitem{CB} R. Cushman and L. Bates,
\emph{Global Aspects of Classical Integrable Systems}, 2nd ed.,
Birkh\"auser, 2015. DOI: \href{https://doi.org/10.1007/978-3-0348-0918-4}
{10.1007/978-3-0348-0918-4}.
\bibitem{Audin} M. Audin, \emph{Spinning Tops: A Course on Integrable
Systems}, Cambridge University Press, 1996, ISBN 978-0-521-56129-7.
\end{thebibliography}
\fi
\end{document}
